\section{Angular hierarchy and predicted factors}
\label{hier:section}

We continue with the finite popular forest and the set $Z_0$ from
\Cref{count:simultaneous}, on which the top estimates, the bound
\eqref{eq:src-12}, and the separated-top count hold.  The parameters
$p,G,K,B$ are defined in \eqref{forest:eq-parameters}, and
$\alpha,\mu,r_*$ in \eqref{phase:parameters}.
We use $i$ for individual angular levels and $j$ for blocks of levels.
Each shift chooses different block endpoints.  The analytic family
constructed below is shared by all $K$ groupings.  Its block identities
explain how the phase estimate will compare a parent amplitude with its
children, while its equal-coordinate value leaves the boundary weights
removed in the final angular-grid average.  We then prove the estimates
that predict amplitudes at a finer angular level, up to a small error,
from one good point on a much longer horizontal interval.  No mass cutoff
is used in those estimates.

All sums below are finite.  We retain the strictly positive spatial
weights $h_T$ from \eqref{eq:src-5}; in particular,
\begin{equation}
 0<h_T\le C_3,
 \qquad
 \int_{\R^2}h_T\,\dd z\le L^C |R_T|.
 \label{hier:top-weight-basic}
\end{equation}
The integral follows by the change of variables to the normalized
coordinates of $R_T$, since the decay exponent in that definition is
greater than two.  Its constant includes the fixed polynomial dilation
$K_0^2$.  Constants denoted by $L^C$ in this section are independent of the
number of tiles, tops, levels, and Hilbert-space components.

\subsection{Eligible masses and shifted blocks}

Choose a dyadic number $d(0)$ such that $B^2\le d(0)<2B^2$, enlarging it
by a fixed absolute factor if necessary to cover the initial precision
range.  For a translation parameter $\vartheta\in[0,d(0))$, let
\begin{equation}
 d(i)=d(0)G^{-i},
 \qquad
 \mathcal G_i^\vartheta
 =\bigl\{[\vartheta+q d(i),\vartheta+(q+1)d(i)):q\in\Z\bigr\}.
 \label{hier:angular-grid}
\end{equation}
Since $G$ is an integer dyadic power, these grids are nested.  We fix
$\vartheta$ for now.  The top slopes lie in a bounded interval, so only an
absolute number of intervals in $\mathcal G_0^\vartheta$ contain a top
slope.  These intervals are the roots.  The half-open convention assigns
each top $T$ a unique interval $v_i(T)$ at every level $i$.

For $v\in\mathcal G_i^\vartheta$ put
\begin{equation}
 \mathcal T(v)
 =\{T:s_T\in v,\ d_T\le d(i)/B^2\},
 \qquad
 n_v(z)=\sum_{T\in\mathcal T(v)}h_T(z).
 \label{hier:eligible-mass}
\end{equation}
We call these tops eligible at $v$.  A node is nonempty if
$\mathcal T(v)\ne\varnothing$.  For such a node $n_v(z)>0$ at every
spatial point.  An empty node has mass and amplitude zero and is omitted
from every expression involving a negative power or logarithm of its
mass.  Thus none of the logarithms below requires a spatial lower cutoff.

If $v_1,\ldots,v_G$ are the immediate children of $b$, then their eligible
top sets are disjoint subsets of $\mathcal T(b)$.  More precisely,
\begin{equation}
 n_b-\sum_{v\text{ child of }b}n_v
 =\sum_{\substack{T\in\mathcal T(b)\\d_T>d(i+1)/B^2}}h_T\ge0,
 \qquad b\in\mathcal G_i^\vartheta.
 \label{hier:mass-loss}
\end{equation}
The same assertion holds between any two levels, with the corresponding
eligibility cutoff.  In particular, every positive mass difference used
below is itself a sum of the same positive functions $h_T$.

For a packet $s$ and a top $T$, define their last eligible levels by
\begin{equation}
 k_s=\max\{i\ge0:d_s\le d(i)/B^2\},
 \qquad
 k_T=\max\{i\ge0:d_T\le d(i)/B^2\}.
 \label{hier:stopping-levels}
\end{equation}
The initial precision choice makes these sets nonempty.  The assigned-top
relation $\ell_s\le l_{T(s)}$ gives $d_{T(s)}\le d_s$, and therefore
$k_s\le k_{T(s)}$.  Thus a packet is stopped only while its assigned top
is still eligible.  Since the collection is finite and all precisions are
positive, the maximum of the $k_T$ is finite.  The grids may be defined at
every level, but no amplitude requires a level beyond this maximum.  We
also keep the first block endpoint past this maximum, so that terminal
blocks have the same notation as other blocks.

For each $\sigma\in\{0,\ldots,K-1\}$ use the endpoints
\begin{equation}
 i_0^\sigma=0,
 \qquad i_j^\sigma=1+\sigma+(j-1)K\quad(j\ge1).
 \label{hier:shifted-blocks}
\end{equation}
The first block has between one and $K$ transitions; all subsequent blocks
have $K$.  The transition $i\to i+1$ is the final transition of a block
for precisely the shift $\sigma\equiv i\pmod K$.  For this shift, let
$a(i)$ be the starting level of that block.  Given an immediate child
$v\in\mathcal G_{i+1}^\vartheta$ with immediate parent
$b\in\mathcal G_i^\vartheta$, denote its ancestor at level $a(i)$ by
$J(i,v)$.  Define
\begin{equation}
 m_v=n_b+
 \sum_{\substack{v'\in\mathcal G_{i+1}^\vartheta,
                    \ v'\subset J(i,v),\ v'\not\subset b\\
                    \dist(v,v')<G^{-1/4}d(i)}}n_{v'}.
 \label{eq:src-14}
\end{equation}
Here the distance is the distance between the angular intervals.  The
summands outside $b$ have disjoint top sets, and these sets are disjoint
from $\mathcal T(b)$.  They all consist of tops eligible at $J(i,v)$.
Consequently,
\begin{equation}
 n_v\le n_b\le m_v\le n_{J(i,v)}.
 \label{hier:boundary-mass-bounds}
\end{equation}
In particular, $m_v-n_v$ is a sum of positive top weights.  The use of the
larger block ancestor in \eqref{eq:src-14} never changes the radius of
the neighbor test: that radius is always tied to the immediate transition
scale $d(i)$.
\Cref{fig:shifted-blocks} shows the shifted block endpoints and the top
sets entering $m_v$.

\begin{figure}[tbp]
\centering
\begin{tikzpicture}[x=.94cm,y=.74cm,>=Stealth,font=\small]
 \foreach \i in {0,...,9}
   \node[font=\scriptsize] at (\i,.62) {\(\i\)};
 \node[anchor=east,font=\scriptsize] at (-.55,.62) {level};
 \foreach \sig/\col in {0/blue!65!black,1/green!45!black,2/orange!80!black}{
   \pgfmathsetmacro{\yy}{-1.18*\sig}
   \node[anchor=east] at (-.5,\yy) {\(\sigma=\sig\)};
   \draw[gray!50] (0,\yy)--(9,\yy);
   \foreach \i in {0,...,8}{
     \pgfmathtruncatemacro{\rr}{mod(\i,3)}
     \ifnum\rr=\sig\relax
       \draw[->,very thick,\col] (\i+.1,\yy)--(\i+.89,\yy);
     \else
       \draw[->,gray!45] (\i+.14,\yy)--(\i+.86,\yy);
     \fi
   }
   \foreach \i in {0,...,9}
     \fill[gray!65] (\i,\yy) circle (1.35pt);
   \pgfmathtruncatemacro{\first}{1+\sig}
   \draw[\col,thick] (0,\yy+.13)--(0,\yy+.32)
     --(\first,\yy+.32)--(\first,\yy+.13);
   \foreach \b in {0,1}{
     \pgfmathtruncatemacro{\aa}{\first+3*\b}
     \pgfmathtruncatemacro{\bb}{\aa+3}
     \draw[\col,thick] (\aa,\yy+.13)--(\aa,\yy+.32)
       --(\bb,\yy+.32)--(\bb,\yy+.13);
   }
 }
\end{tikzpicture}

\medskip
\begin{tikzpicture}[x=.94cm,y=1cm,font=\small,>=Stealth]
 \draw[thick] (0,0)--(9,0);
 \draw[thick] (0,-.12)--(0,.12) (9,-.12)--(9,.12);
 \node[above] at (4.5,.13) {block ancestor \(J\)};
 \fill[blue!12] (2.3,-1.35) rectangle (6,-.7);
 \draw[blue!65!black,thick] (2.3,-1.35) rectangle (6,-.7);
 \node at (3.85,-1.025) {immediate parent \(b\)};
 \fill[blue!65!black] (5.48,-1.35) rectangle (5.86,-.7);
 \node[below] at (5.67,-1.4) {\(v\)};
 \foreach \aa in {6.06,6.48}{
   \fill[orange!35] (\aa,-1.35) rectangle (\aa+.34,-.7);
   \draw[orange!80!black] (\aa,-1.35) rectangle (\aa+.34,-.7);
 }
 \node[anchor=west] at (7.15,-1.02) {neighbors \(v'\)};
 \draw[thin] (6.84,-1.02)--(7.08,-1.02);
 \draw[dashed] (5.48,-1.7)--(5.48,-2.4);
 \draw[dashed] (6.95,-1.5)--(6.95,-2.4);
 \draw[<->] (5.48,-2.27)--(6.95,-2.27);
 \node[anchor=west] at (7.05,-2.27) {\(\sim G^{-1/4}d(i)\)};
 \node[anchor=east] at (4.9,-2.27) {local neighbor radius};
\end{tikzpicture}
\caption{Shifted blocks and boundary masses, schematically.
The upper diagram uses \(K=3\). Each shift starts with \(1+\sigma\)
transitions and then uses blocks of \(K\) transitions; colored arrows
are the final transitions for that shift.
At a final transition \(i\to i+1\), the mass \(m_v\) includes all of
\(n_b\) and the masses of neighboring children outside \(b\) but inside
\(J\). The neighbor radius is determined by the immediate scale \(d(i)\),
even when \(J\) begins much earlier. Neither diagram is to scale.}
\label{fig:shifted-blocks}
\end{figure}


For $\zeta=(\zeta_0,\ldots,\zeta_{K-1})\in\C^K$ and nonempty $v$ at
level $i+1$, put
\begin{equation}
 w_v(\zeta,z)
 =\exp\!\left(\zeta_{i\bmod K}\log\frac{n_v(z)}{m_v(z)}\right).
 \label{hier:analytic-transition}
\end{equation}
All logarithms in this formula are real logarithms of positive numbers.
Thus the weights are entire functions of $\zeta$ and have modulus at
most one when all $\Re\zeta_\sigma\ge0$.

\subsection{The analytic family and its block comparisons}

For a nonempty node $v$ at level $h$, define the $\mathcal H$-valued
analytic amplitude by
\begin{equation}
 A_v(\zeta,z)
 =\sum_{\substack{s:s_{T(s)}\in v\\ k_s\ge h}}
 q_{T(s)}(z)\phi_s(z)\,
 n_{v_{k_s}(T(s))}(z)^{-\mu/K}
 \prod_{i=h}^{k_s-1}w_{v_{i+1}(T(s))}(\zeta,z).
 \label{eq:src-15}
\end{equation}
The condition $k_s\ge h$ already ensures eligibility of $T(s)$ at $v$.
The empty product is one.  This single definition is used for every
shift.  At fixed $z$, it is a finite sum of entire $\mathcal H$-valued
functions of $\zeta$.

We first exhibit two algebraic uses of this common family.  For the block
comparison, temporarily fix a shift $\sigma$ and parameters satisfying
\begin{equation}
 \Re\zeta_\sigma=\mu,
 \qquad \Re\zeta_{\sigma'}=0\quad(\sigma'\ne\sigma),
 \qquad \norm{\Im\zeta}_2\le L^2.
 \label{hier:vertex-parameters}
\end{equation}
We refer to this choice of real parts as the vertex corresponding to
$\sigma$.  Write
$i_j=i_j^\sigma$ and $\rho_j=d(i_j)$.  Thus
\begin{equation}
 G\le\frac{\rho_j}{\rho_{j+1}}\le B,
 \label{hier:block-ratio}
\end{equation}
including the first, possibly shorter, block.

For a nonempty $J\in\mathcal G_{i_j}^\vartheta$, call its nonempty
descendants at level $i_{j+1}$ its \emph{block children}.  For such a block child $c$,
all packets entering $A_c$ have the same transition factor from $J$ to
$c$.  Denote it by
$\lambda_c=\prod_{i=i_j}^{i_{j+1}-1}w_{v_{i+1}}$, along that descendant
path.  Only the last transition of this block has a nonzero real
exponent.  Separating packets according to whether their stopping level
precedes $i_{j+1}$ therefore gives the exact identity
\begin{equation}
 A_J=\sum_{\substack{c\subset J\\c\in\mathcal G_{i_{j+1}}^\vartheta}}
                 \lambda_c A_c+B_J,
 \qquad
 |\lambda_c|=\left(\frac{n_c}{m_c}\right)^\mu.
 \label{eq:src-18}
\end{equation}
The remainder $B_J$ is precisely the sum in \eqref{eq:src-15} over
$i_j\le k_s<i_{j+1}$.  Its assigned tops need not lose eligibility at
the end of the block.  Thus residual packets and lost top mass are
different parts of the block decomposition.

For each nonempty $J$, we add the mass to the normalized $p$-th power
of the amplitude.  The resulting positive amount $V_J$ and its
multiplicative comparison factor $\beta_J$ are
\begin{equation}
 V_J=n_J+n_J^{\alpha-p}\norm{A_J}_{\mathcal H}^p,
 \qquad
 \beta_J=
 \frac{V_J}{n_J+\sum_c n_c^{\alpha-p}\norm{A_c}_{\mathcal H}^p}.
 \label{eq:src-19}
\end{equation}
We use $\beta_J^0$ for the same ratio with $B_J$ omitted only in the
numerator.  The mass baseline makes $V_J$ and every denominator positive,
including when all its block-child amplitudes vanish.  The sum is over
the nonempty descendants at the end of the block.  Equivalently, the
denominator is the lost eligibility mass $n_J-\sum_c n_c$ plus the
block-child amounts $\sum_c V_c$.  This is the comparison later used
in the exact path expansion.

At any one spatial point, abbreviate $n=n_J$ and normalize the child
amplitudes by
\begin{equation}
 l_c=\frac{n_c}{n},\qquad
 X_c=n_c^{\mu-1}A_c,\qquad
 \beta_J^0=
  \frac{1+\norm{\sum_c l_c^{1-\mu}\lambda_c X_c}_{\mathcal H}^p}
       {1+\sum_c l_c\norm{X_c}_{\mathcal H}^p}.
 \label{eq:src-25}
\end{equation}
The ratio follows by dividing \eqref{eq:src-19} by $n$ and using
$p\mu=\alpha-1$.  Also $\sum_c l_c\le1$.  This is a pointwise identity:
the masses, scalar factors, and vectors still vary with the spatial
point.  The vector sum in its numerator is the weighted child sum whose
phase motion we will compare with the weighted child energies in the
denominator.

The boundary mass $m_c$ is designed to pay for selecting separated
children in that comparison.  At the final transition of this block put
$i=i_{j+1}-1$ and
\begin{equation}
 Q=G^{-1/4}d(i)=G^{3/4}\rho_{j+1}.
 \label{hier:child-separation}
\end{equation}
Fix one spatial point and choose any finite subfamily $\mathcal C$ of
nonempty block children.  Join distinct $c,c'\in\mathcal C$ when
$\dist(c,c')<Q$.  Give the children independent exponential clocks of
rates $n_c$, and retain a child when its clock rings before those of all
its neighbors.  The rates are positive.  The clocks have no ties with
probability one, and the retained intervals have pairwise distance at
least $Q$.  The retention probability of $c$ is
\begin{equation}
 \pi_c=\int_0^\infty n_c e^{-n_ct}
              \prod_{c'\sim c}e^{-n_{c'}t}\,\dd t
       =\frac{n_c}{n_c+\sum_{c'\sim c}n_{c'}}
       \ge\frac{n_c}{m_c}.
 \label{pred:retention-probability}
\end{equation}
To see the inequality, let $b$ be the immediate parent of $c$.  The
eligible top sets of $c$ and its competitors inside $b$ are disjoint
subsets of $\mathcal T(b)$, so their masses sum to at most $n_b$.
Every competitor outside $b$ lies inside $J$ and within
$Q=G^{-1/4}d(i)$ of $c$, so its mass is included among the additional
terms of $m_c$ in \eqref{eq:src-14}.  Thus $m_c$ bounds the entire
competing mass even when only a subfamily of children is used.

Multiplying the coefficient of a retained child by $\pi_c^{-1}$ reconstructs its original
coefficient in expectation.  The corresponding expected $r_*$-power
cost is
\[
 \E\!\left[\one_{\{c\text{ retained}\}}\pi_c^{-1}\right]=1,
 \qquad
 \E\!\left[\one_{\{c\text{ retained}\}}\pi_c^{-r_*}\right]
       =\pi_c^{1-r_*}.
\]
The coefficient exponent in \eqref{eq:src-18} matches this cost
exactly.  The identities $r_*(1-\mu)=1$ and $\mu r_*=r_*-1$ give
\begin{equation}
 \pi_c^{1-r_*}|\lambda_c|^{r_*}
   =\left(\frac{n_c/m_c}{\pi_c}\right)^{r_*-1}\le1.
 \label{hier:selection-cancellation}
\end{equation}
The first identity also turns $l_c^{r_*(1-\mu)}$ into $l_c$, the weight
in the denominator of \eqref{eq:src-25}.  The later snapshot construction
will use this selection for the child vectors it retains and check the
remaining hypotheses of the finite Fourier estimate.

The same analytic family has a second algebraic use at the equal-coordinate
parameter
\[
 \zeta^{\mathrm c}=(\mu/K,\ldots,\mu/K).
\]
Let $v$ be a root and let $s$ be a packet contributing to $A_v$, with
$T=T(s)$.  Apart from the factor $q_T\phi_s$, the coefficient of $s$ in
$n_v^{\mu/K}A_v(\zeta^{\mathrm c},z)$ is
\[
 n_v^{\mu/K}n_{v_{k_s}(T)}^{-\mu/K}
       \prod_{i=0}^{k_s-1}
        \left(\frac{n_{v_{i+1}(T)}}{m_{v_{i+1}(T)}}\right)^{\mu/K}.
\]
Cancelling consecutive eligible masses identifies it as
\begin{equation}\label{eq:src-30}
 \kappa_s
   =\prod_{i=0}^{k_s-1}
       \left(\frac{n_{v_i(T)}}{m_{v_{i+1}(T)}}\right)^{\mu/K},
 \qquad T=T(s).
\end{equation}
The empty product is one, and each factor lies in $(0,1]$ by
\eqref{eq:src-14}.  Thus for each root the equal-coordinate value produces
its portion of the masked packet sum with only these boundary ratios in
its coefficients.  The
remaining tasks are to estimate the vertex values, interpolate those
estimates to $\zeta^{\mathrm c}$, and remove the boundary ratios by
averaging the angular-grid translation.

\subsection{Full-root variation and amplitude bounds}

We now leave the temporary vertex specialization and return to the general
range $0\le\Re\zeta_\sigma\le\mu$ for every coordinate, with arbitrary
imaginary parts.  The following bounds control the endpoint weights
without a loss proportional to the number of levels.  Along one eligible
top path, write $n_i=n_{v_i(T)}$ and
\begin{equation}
 P_i=\log\frac{n_i}{n_{i+1}},
 \qquad
 R_i=\log\frac{m_{v_{i+1}(T)}}{n_{i+1}}.
 \label{hier:positive-logs}
\end{equation}
Both are nonnegative as long as the top remains eligible at level $i+1$.

\begin{lemma}[Full-root variation bound]
\label{hier:weight-variation}
Fix an eligible path from a root $v_0(T)$ through level $k$.  At every
spatial point,
\begin{equation}
 \sum_{i=0}^{k-1}R_i
 \le K\log\frac{n_{v_0(T)}}{n_{v_k(T)}},
 \qquad
 \sum_{i=0}^{k-1}P_i
 =\log\frac{n_{v_0(T)}}{n_{v_k(T)}}.
 \label{hier:full-root-telescope}
\end{equation}
For a fixed starting level $h\le k_T$, define, for $h\le k\le k_T$,
\[
 \omega_{T,h,k}(\zeta,z)
 =n_{v_k(T)}(z)^{-\mu/K}
   \prod_{i=h}^{k-1}w_{v_{i+1}(T)}(\zeta,z).
\]
If $0\le\Re\zeta_\sigma\le\mu$ for every $\sigma$, then the supremum
plus total variation of this sequence, with the endpoint $k$ varying,
is bounded by
\begin{equation}
 \sup_{h\le k\le k_T}|\omega_{T,h,k}|
 +\sum_{k=h}^{k_T-1}
       |\omega_{T,h,k+1}-\omega_{T,h,k}|
 \le C h_T^{-\mu/K}
 \left[1+K(1+\norm{\zeta}_2)
                   \log\frac{n_{v_0(T)}}{h_T}\right].
 \label{eq:src-16}
\end{equation}
The same bound holds upon restricting the endpoints to a subinterval.
\end{lemma}

\begin{proof}
Fix a shift $\sigma$.  For the transitions $i\equiv\sigma\pmod K$
with $i<k$, \eqref{hier:boundary-mass-bounds} gives
\[
 R_i\le\log\frac{n_{a(i)}}{n_{i+1}}.
\]
For this fixed shift, the intervals of levels $[a(i),i+1]$ are successive
blocks starting at zero.  Their logarithms telescope.  The endpoint of
the last completed block is at most $k$, so the sum is at most
$\log(n_0/n_k)$, by monotonicity of the masses.  Summing over the $K$
shifts proves the first inequality in
\eqref{hier:full-root-telescope}.  The second is an ordinary telescope.
This is a full-root bound.  In particular, it remains applicable to
weights occurring on a later segment even when another shift's block
starts before that segment.

Set $a=\mu/K$, solely for this proof, and use the specified logarithm
\[
 L_k=-a\log n_k-\sum_{i=h}^{k-1}\zeta_{i\bmod K}R_i,
 \qquad \omega_{T,h,k}=e^{L_k}.
\]
Its increment is
\begin{equation}
 L_{k+1}-L_k=aP_k-\zeta_{k\bmod K}R_k.
 \label{hier:log-weight-increment}
\end{equation}
Since $n_k\ge h_T$ and all real transition exponents are nonnegative,
$|e^{L_k}|\le h_T^{-a}$.  For any complex $A,B$,
\[
 |e^B-e^A|
 \le |B-A|\int_0^1|e^{(1-t)A+tB}|\,\dd t
 \le |B-A|\max(|e^A|,|e^B|).
\]
Using this inequality for adjacent logarithms and then
\eqref{hier:full-root-telescope} proves \eqref{eq:src-16}.  Notice that
the exponential of the sum of the absolute logarithmic increments has
not been used.  Such an exponential would be unsuitable when $h_T$ is
very small.
\end{proof}

\begin{lemma}[Amplitude bounds]
\label{hier:amplitudes}
On $Z_0$, for $0\le\Re\zeta_\sigma\le\mu$,
\begin{equation}
 \norm{A_v(\zeta,z)}_{\mathcal H}
 \le L^C(1+\norm{\zeta}_2)n_v(z).
 \label{eq:src-17}
\end{equation}
More precisely, the contribution of any one eligible top $T$ to $A_v$,
also when restricted to any interval of packet lengths, is bounded by
\begin{equation}
 L^C(1+\norm{\zeta}_2)|q_T(z)|h_T(z)^4.
 \label{hier:individual-amplitude}
\end{equation}
If all eligible tops of $v$ are insignificant at $z\in Z_0$, then
\begin{equation}
 \norm{A_v(\zeta,z)}_{\mathcal H}
 \le L^C(1+\norm{\zeta}_2)e^{-L\log L}n_v(z).
 \label{hier:insignificant-amplitude}
\end{equation}
\end{lemma}

\begin{proof}
For vectors $f_h,\ldots,f_k$ and scalars $\omega_h,\ldots,\omega_k$,
summation by parts gives
\begin{equation}
 \norm{\sum_{i=h}^k\omega_i f_i}
 \le
 \left(|\omega_k|+\sum_{i=h}^{k-1}|\omega_{i+1}-\omega_i|\right)
 \max_{h\le j\le k}\norm{\sum_{i=h}^j f_i}.
 \label{hier:summation-by-parts}
\end{equation}
Group the packets assigned to $T$ according to their stopping level.
The resulting initial sums, or the sums obtained after restricting to an
interval of packet lengths, are interval partial sums of the original
top's length sequence.  By \eqref{eq:src-11} and the fact that the
$K_0$-dilated weight is at least the corresponding $K_1$-dilated weight,
their norm at $z\in Z_0$ is at most $L^C h_T(z)^8$.  The mask $q_T(z)$
is common to this whole top.  Apply
\eqref{hier:summation-by-parts} and \eqref{eq:src-16}.

To see that the logarithmic factor can be absorbed uniformly, use
\eqref{eq:src-12} and $\log N\le2L$.  Since $n_{v_0(T)}\le n$ and
$h_T\le C_3$,
\[
 0\le\log\frac{n_{v_0(T)}}{h_T}
 \le C L+\log^+\frac1{h_T}.
\]
Here and below $\log^+x=\max(0,\log x)$.  Also $\mu/K<1$.  The
elementary boundedness of $x^\delta\log^q(1/x)$ on $(0,1]$, for fixed
$\delta>0$ and fixed $q$, shows that
\[
 h_T^{8-\mu/K}
 \left[1+K(1+\norm{\zeta}_2)
                    \log\frac{n_{v_0(T)}}{h_T}\right]
 \le L^C(1+\norm{\zeta}_2)h_T^4.
\]
This proves \eqref{hier:individual-amplitude}.  Summing over the eligible
tops and using $h_T^4\le C_3^3h_T$ proves \eqref{eq:src-17}.
For insignificant tops, the mask estimate after \eqref{eq:src-8} gives
$|q_T|\le e^{-L\log L}$.  Applying the same summation proves
\eqref{hier:insignificant-amplitude}.
\end{proof}

\subsection{Spatial meshes for prediction}

Return to a fixed shift $\sigma$ and a vertex parameter satisfying
\eqref{hier:vertex-parameters}.  Write $i_j=i_j^\sigma$ and
$\rho_j=d(i_j)$ as in the block comparison above, and take
$J\in\mathcal G_{i_j}^\vartheta$.  Its block children and their factors
are those in \eqref{eq:src-18}.
Let $R$ be any top eligible in $J$, and parametrize a line of slope
$s_R$ by $z(x)=(x,y_*+s_Rx)$.  On this line use one nested dyadic grid in
the parameter $x$, at lengths
\begin{equation}
 H_j=\frac{w}{\rho_{j+1}\sqrt G},
 \qquad H_{-1}=\frac{w}{\rho_0\sqrt G}.
 \label{eq:src-20}
\end{equation}
These are dyadic lengths because $w,d(0),G$, and $\sqrt G$ are dyadic.
They increase with $j$.  The grids can, for example, have common origin
zero.  Each $H_j$ cell is a union of $H_{j-1}$ cells.

\subsection{Positive-mass motion and the locality of other shifts}

\begin{lemma}[Mass motion]
\label{hier:mass-motion}
Fix $R,J,j$ as above and an interval $I$ of horizontal length $H_j$ on a
line of slope $s_R$.  Consider the masses $n_J,n_c,m_c$ in
\eqref{eq:src-18}, the masses used in the factors $\lambda_c$, and those
used in any endpoint multiplier of a packet contributing to $A_c$.
Each such mass, and each of its positive subset-complement differences
used below, is a sum of weights of tops $T$ satisfying
\begin{equation}
 |s_T-s_R|\le C\rho_j,
 \qquad d_T\le\rho_j/B^2.
 \label{hier:motion-locality}
\end{equation}
There is an absolute constant $C$ such that, with
$\varepsilon_{\mathrm{mov}}=C G^{-1/2}$, every nonzero such mass $a$ satisfies
\begin{equation}
 |\partial_x\log a|\le\frac{\varepsilon_{\mathrm{mov}}}{H_j}
 \quad\text{almost everywhere on }I.
 \label{hier:mass-log-derivative}
\end{equation}
If $r=a/(a+b)$ is one of the ratios of a submass to its containing mass,
then
\begin{equation}
 |\partial_x\log r|
 \le\frac{2\varepsilon_{\mathrm{mov}}}{H_j}(1-r)
 \le\frac{2\varepsilon_{\mathrm{mov}}}{H_j}(-\log r).
 \label{hier:ratio-derivative}
\end{equation}
When $b=0$, the ratio is identically one and its derivative is zero.
In particular, nonzero $n_J-n_c$ and $m_c-n_c$ have the relative motion
bound \eqref{hier:mass-log-derivative}.  All these masses vary by a
factor $1+O(G^{-1/2})$ across $I$.

For the parameters \eqref{hier:vertex-parameters}, the block
factors also satisfy, for $x,x_0\in I$,
\begin{equation}
 \left|\frac{\lambda_c(x)}{\lambda_c(x_0)}-1\right|
 \le G^{-c_0}
 \label{hier:complex-transition-motion}
\end{equation}
for some absolute $c_0>0$ and all sufficiently large $L$.
\end{lemma}

\begin{proof}
All descendants of $J$ have slopes inside $J$, so their top slopes are
within $\rho_j$ of $s_R$.  Their eligibility cutoffs are at most
$\rho_j/B^2$.  A mass in an additional boundary term
\eqref{eq:src-14} at a transition $i\ge i_j$ can come from an interval
outside $J$, because the block associated to another shift can start
before $i_j$.  However, that interval is within
$G^{-1/4}d(i)$ of the relevant child and has width $d(i+1)$.
The relevant child is a descendant of $J$, and $d(i)\le\rho_j$.
Hence every added slope is within $C\rho_j$ of $s_R$.  The tops in an
added interval have precision at most $d(i+1)/B^2$; those in the
immediate-parent mass have precision at most $d(i)/B^2$.
This proves \eqref{hier:motion-locality}, including for weights below
the block.  Taking a difference of the relevant nested top sets preserves
the same property.  This argument uses the neighbor radius at the
immediate transition, rather than the possibly larger width of the other
shift's block ancestor.

In normalized coordinates of a top $T$, horizontal motion along $s_R$
has velocities $1/l_T$ and $(s_R-s_T)/W$.  Differentiating
\eqref{eq:src-5}, using the Lipschitz property of the absolute value at
almost every point, gives
\[
 |\partial_x\log h_T|
 \le C\left(\frac1{l_T}+\frac{|s_R-s_T|}{W}\right)
 \le C\left(\frac{\rho_j}{B^2w}+\frac{\rho_j}{mw}\right).
\]
The factor $K_0^{-1}$ available from the dilation is not needed.
Multiplying by $H_j$ and using \eqref{hier:block-ratio} yields
\[
 H_j|\partial_x\log h_T|
 \le\frac C{\sqrt G}
       \left(\frac{B}{B^2}+\frac Bm\right)
 \le C G^{-1/2}.
\]
For a positive sum $a=\sum_T h_T$, its logarithmic derivative is the
weighted average of the logarithmic derivatives of the summands.
This proves \eqref{hier:mass-log-derivative}, also for positive
differences, which are such sums by their set interpretation.

If $a,b>0$, direct differentiation gives
\[
 \partial_x\log\frac a{a+b}
 =\frac b{a+b}\bigl(\partial_x\log a-\partial_x\log b\bigr).
\]
This proves \eqref{hier:ratio-derivative}; the inequality
$1-r\le-\log r$ holds for $0<r\le1$.  If the complementary top set is
empty, $r=1$ everywhere.  In particular, writing $q_r=-\log r\ge0$,
\begin{equation}
 |\partial_x q_r|\le\frac{2\varepsilon_{\mathrm{mov}}}{H_j}q_r,
 \qquad
 e^{-2\varepsilon_{\mathrm{mov}}}q_r(x_0)
 \le q_r(x)\le e^{2\varepsilon_{\mathrm{mov}}}q_r(x_0)
 \quad(x,x_0\in I).
 \label{hier:loss-motion}
\end{equation}
The inequalities remain valid for an identically zero loss.

Finally, a block factor is a product of at most $K$ transition factors.
The simpler bound $|\partial_x\log r|\le2\varepsilon_{\mathrm{mov}}/H_j$ shows that
its specified complex logarithm changes by at most
$C K(1+\norm{\zeta}_2)G^{-1/2}$.  For
\eqref{hier:vertex-parameters} this is $L^C G^{-1/2}$ and tends to zero
faster than any negative fixed power of $L$.  Exponentiating proves
\eqref{hier:complex-transition-motion} after decreasing $c_0$.
This last argument does not require any bound on the size of a
logarithmic loss itself.
\end{proof}

\subsection{Variation of moving endpoint weights}

The next estimate makes explicit why the motion of the analytic weights
can be controlled from a single point, even for very small $h_T$.
For a finite scalar sequence $a=(a_k)$ let
\[
 \norm{a}_{\mathrm{BV}}
 =\sup_k|a_k|+\sum_k|a_{k+1}-a_k|.
\]
Only this elementary finite-sequence norm is used here.

\begin{lemma}[Moving weights]
\label{hier:moving-weights}
In the setting of \Cref{hier:mass-motion}, fix a top $T$ contributing to
$A_c$, and take endpoint weights $\omega_k=\omega_{T,i_{j+1},k}$.
If $z_0\in I\cap Z_0$ and the parameters satisfy
\eqref{hier:vertex-parameters}, then, uniformly for $z'\in I$,
\begin{align}
 \norm{(\omega_k(z'))_k}_{\mathrm{BV}}
 &\le L^C h_T(z_0)^{-4},
 \label{hier:moving-weight-bv}\\
 \norm{(\omega_k(z')-\omega_k(z_0))_k}_{\mathrm{BV}}
 &\le G^{-1/2}L^C h_T(z_0)^{-4}.
 \label{hier:moving-weight-difference}
\end{align}
Both conclusions hold for any interval of endpoints.  Neither assumes
that $z'$ belongs to $Z_0$.
\end{lemma}

\begin{proof}
Parametrize the interval by its horizontal coordinate $x$, and write
$h=i_{j+1}$ and $a=\mu/K$ in this proof.  Use $P_i,R_i$ from
\eqref{hier:positive-logs} for $h\le i<k_T$.  By
\eqref{hier:loss-motion}, each of these nonnegative losses is comparable
throughout $I$ to its value at $x_0$, with factor
$e^{O(\varepsilon_{\mathrm{mov}})}$, and its absolute derivative is bounded by
$C\varepsilon_{\mathrm{mov}}/H_j$ times itself.  Define
\[
 \mathcal B(x)=\sum_{i=h}^{k_T-1}
                    \bigl(aP_i(x)+|\zeta_{i\bmod K}|R_i(x)\bigr).
\]
It follows that $\mathcal B(x)\le e^{C\varepsilon_{\mathrm{mov}}}\mathcal B(x_0)$.
By the full-root bound,
\begin{equation}
 \mathcal B(x_0)
 \le C K(1+\norm{\zeta}_2)
                 \log\frac{n_{v_0(T)}(z_0)}{h_T(z_0)}.
 \label{hier:moving-loss-budget}
\end{equation}
The full-root mass is used only at the snapshot in this inequality; its
spatial derivative is not needed.

Let $L_k=\log\omega_k$ be the specified logarithm in the proof of
\Cref{hier:weight-variation}, and put $q_k=\partial_x L_k$.
Equation~\eqref{hier:log-weight-increment} and the mass derivative
bound imply
\begin{align*}
 \sup_k|q_k|
 &\le\frac{C\varepsilon_{\mathrm{mov}}}{H_j}
                  \bigl(a+\mathcal B(x)\bigr),\\
 \sum_k|q_{k+1}-q_k|
 &\le\frac{C\varepsilon_{\mathrm{mov}}}{H_j}\mathcal B(x).
\end{align*}
For example, the initial term is
$q_h=-a\partial_x\log n_h$, and subsequent differences are the
derivatives of $aP_i-\zeta_{i\bmod K}R_i$.  At every point of $I$,
\[
 \sup_k|\omega_k(x)|\le h_T(x)^{-a}
 \le e^{a\varepsilon_{\mathrm{mov}}}h_T(z_0)^{-a}.
\]
The complex-exponential estimate used for
\eqref{eq:src-16} therefore gives
\[
 \norm{\omega(x)}_{\mathrm{BV}}
 \le C h_T(z_0)^{-a}\bigl(1+\mathcal B(x_0)\bigr).
\]

For finite sequences the product rule for differences gives
\[
 \norm{ab}_{\mathrm{BV}}
 \le \sup_k|a_k|\norm{b}_{\mathrm{BV}}
       +\sup_k|b_k|\sum_k|a_{k+1}-a_k|.
\]
Using $\partial_x\omega_k=\omega_k q_k$ and the preceding estimates
yields
\begin{equation}
 \norm{\partial_x\omega(x)}_{\mathrm{BV}}
 \le\frac{C\varepsilon_{\mathrm{mov}}}{H_j}
 h_T(z_0)^{-a}\bigl(1+\mathcal B(x_0)\bigr)^2.
 \label{hier:weight-derivative-bv}
\end{equation}
The functions involved are absolutely continuous on the line.  Integrate
\eqref{hier:weight-derivative-bv} between $x_0$ and $x'$, over distance
at most $H_j$, and use the triangle inequality for the finite-sequence
norm.  This proves the difference bound with right side
$C\varepsilon_{\mathrm{mov}} h_T(z_0)^{-a}(1+\mathcal B(x_0))^2$.

At the snapshot, \eqref{eq:src-12},
\eqref{hier:vertex-parameters}, and
\eqref{hier:moving-loss-budget} bound the last expression by
$G^{-1/2}L^C h_T(z_0)^{-4}$.  Indeed the only possible nonpolynomial
quantity is $\log^+(1/h_T(z_0))$, whose square is absorbed by the
difference between exponents $a<1$ and $4$.  The same argument proves
\eqref{hier:moving-weight-bv}.  There is no assumption that
$G^{-1/2}\mathcal B(x_0)$ is small: very small $h_T$ is allowed, and
only a polynomial in its logarithm has appeared.
\end{proof}

\subsection{Amplitude motion from one good snapshot}

\begin{lemma}[Predictable amplitude motion]
\label{hier:amplitude-motion}
Fix a block child $c$ of $J$, a path top $R$ eligible in $J$,
and an $H_j$ cell on a line of slope $s_R$.  Let $t_c$ be the angular
center of $c$.  If $z_0\in Z_0$ lies in that cell and
$z'=z_0+(b,s_Rb)$ is any other point in the cell, then, for the parameters \eqref{hier:vertex-parameters},
\begin{equation}
 \norm{A_c(\zeta,z')-\mathcal D_c(b)A_c(\zeta,z_0)}_{\mathcal H}
 \le G^{-c_0}n_c(z_0),
 \qquad
 \mathcal D_c(b)_{\nu\nu}
 =\exp\bigl(i(s_R-t_c)\eta_\nu b\bigr).
 \label{eq:src-21}
\end{equation}
Here $c_0>0$ is absolute and may also be used, after decreasing it, in
\eqref{hier:complex-transition-motion}.  The estimates hold at all
displaced points in the cell, not just its points in $Z_0$.
\end{lemma}

\begin{proof}
Fix one top $T$ contributing to $A_c$.  Its contributing packets have
$k_s\ge i_{j+1}$, hence
\begin{equation}
 \ell_s\ge l_*:=\frac{B^2w}{\rho_{j+1}},
 \qquad l_T\ge l_*.
 \label{hier:predictable-length}
\end{equation}
Also $s_T\in c$, so $|s_T-t_c|\le\rho_{j+1}/2$, while
$|s_T-s_R|\le C\rho_j$.

Consider any interval partial sum $S$ of these unweighted packets on
$T$.  The good-snapshot estimates \eqref{eq:src-11}, including their
demodulated vertical-derivative and vertical-neighborhood versions, bound
$S$ and $W(\partial_y-i\eta_\nu)S_\nu$ by
$L^Ch_T(z_0)^8$ throughout the vertical neighborhood of radius $W$
at $z_0$.  Put $v=(s_R-s_T)b$ and first move to
$z_1=z_0+(0,v)$.  Since
\[
 |v|\le C\rho_jH_j\le\frac{CBw}{\sqrt G}\ll W,
\]
the fundamental theorem of calculus after componentwise demodulation
gives
\[
 \norm{S(z_1)-\operatorname{diag}(e^{i\eta_\nu v})S(z_0)}
 \le L^Ch_T(z_0)^8\frac{|v|}{W}
 \le G^{-1/2}L^Ch_T(z_0)^8.
\]
Next move from $z_1$ to $z'$ along slope $s_T$ by horizontal distance
$b$.  The nonexceptional horizontal derivative estimate in
\Cref{conc:top-good}, applied to lengths at least $l_*$, bounds this
change by
\[
 L^Ch_T(z_0)^8\frac{|b|}{l_*}
 \le G^{-1/2}L^Ch_T(z_0)^8.
\]
To justify its spatial weight here, the two displacements change the
normalized coordinates of $T$ by at most
$C/(B^2\sqrt G)+CB/(m\sqrt G)$.  Its weight at every intermediate
point is therefore comparable to $h_T(z_0)$.  If the derivative estimate
is expressed using the smaller dilation $K_1$, that weight is bounded
above by the $K_0$-dilated weight used here.  No good-point estimate at
$z_1$ or $z'$ has been invoked for this horizontal move.

The replacement of $s_T$ by $t_c$ in the diagonal phase costs at most
\[
 C|s_T-t_c|\frac{|b|}{w}\norm{S(z_0)}
 \le G^{-1/2}L^Ch_T(z_0)^8,
\]
because $\eta_\nu\asymp w^{-1}$ uniformly in $\nu$.
This proves, uniformly over the interval partial sum,
\begin{equation}
 \norm{S(z')-\mathcal D_c(b)S(z_0)}
 \le G^{-1/2}L^Ch_T(z_0)^8.
 \label{hier:unweighted-motion}
\end{equation}

The same bound holds with the mask $q_T$ multiplying $S$ at its own
evaluation point.  Indeed $|q_T|\le1$, and its formula
\eqref{eq:src-8}, together with boundedness of the derivative of sinc,
gives
\[
 |q_T(z')-q_T(z_0)|
 \le L^C\left(\frac{|b|}{l_T}
                    +\frac{|s_R-s_T|\,|b|}{W}\right)
 \le G^{-1/2}L^C.
\]
Combining this with the snapshot bound for $S$ proves the masked version
of \eqref{hier:unweighted-motion}.  This estimate uses absolute motion
of the mask and does not divide by $q_T$, which may vanish.

Group the masked packets by stopping level and denote their sums by
$f_k(z)$.  Their weighted contribution to the difference in
\eqref{eq:src-21} is
\begin{align*}
 &\sum_k\omega_k(z')
       \bigl(f_k(z')-\mathcal D_c(b)f_k(z_0)\bigr)\\
 &\hspace{15mm}
 +\sum_k\bigl(\omega_k(z')-\omega_k(z_0)\bigr)
                        \mathcal D_c(b)f_k(z_0).
\end{align*}
All interval partial sums of the first parenthesis have norm at most
$G^{-1/2}L^Ch_T(z_0)^8$, by the masked version of
\eqref{hier:unweighted-motion}.  The snapshot partial sums in the second
line have norm at most $L^Ch_T(z_0)^8$.  Apply
\eqref{hier:summation-by-parts} to both lines, then
\eqref{hier:moving-weight-bv} and
\eqref{hier:moving-weight-difference}.  The result for this top is at
most $G^{-1/2}L^Ch_T(z_0)^4$.
Sum over the tops eligible in $c$.  Since $h_T^4\le C_3^3h_T$, this is
at most $G^{-1/2}L^Cn_c(z_0)$.  Finally
$\log G\asymp L^{0.7}$ absorbs every fixed power of $L$ into a smaller
positive power of $G$, proving \eqref{eq:src-21}.
\end{proof}

\subsection{One deterministic residual sequence for each path}

We now refine $Z_0$ by a set which depends on the angular translation but
not on $\zeta$.  Its purpose is to make the later residual estimates
available on most of the path weight.  The subsets of packets used in
its definition must be fixed before evaluating a spatial point or
choosing a mass cutoff.

For each possible path top $R$, let
\begin{equation}
 \mathcal P_R
 =\{s:\ell_s\le l_R,\quad
                    |\theta_s-s_R|\le C B^4d_s\},
 \label{hier:reference-packets}
\end{equation}
where the absolute constant is large enough to cover the assigned-top
comparisons below.  Apply the main/error decomposition of
\Cref{conc:line-moments} with reference slope $s_R$ to this fixed
collection.  Denote the resulting positive error sum by
$\mathcal E_R$ and positive main square function by $\mathcal S_R$.
The absolute envelopes can use fixed decay exponents slow enough to
dominate the norms of the unweighted synthesis packets.  There is no
loss in doing so because the packet bounds are available with arbitrarily
large fixed decay exponents.

For each shift $\sigma$, define a second fixed collection
\begin{equation}
 \mathcal Q_{R,\sigma}
 =\bigcup_{\substack{j\ge0\\ i_{j+1}^\sigma\le k_R}}
 \left\{s:
 \begin{array}{l}
 i_j^\sigma\le k_s<i_{j+1}^\sigma,\\
 T(s)\in\mathcal T\bigl(v_{i_{j+1}^\sigma}(R)\bigr)
 \end{array}\right\}.
 \label{hier:deterministic-residual-subset}
\end{equation}
For blocks at whose endpoint $R$ is still eligible, these are the packets
that stop within the block while their assigned tops remain eligible in
the child $v_{i_{j+1}^\sigma}(R)$.  The union takes all such blocks at
once, fixing one packet collection before any spatial point or mass cutoff
is chosen.  Let $\mathcal V_{R,\sigma}$ be the $3$-variation of the
sequence formed by summing $q_{T(s)}\phi_s$ from this collection at each
packet length.  It permits any finite list of disjoint intervals of that
one length sequence.  Subsequent choices of those intervals may depend
on the evaluation point; the collection and its length sequence remain fixed.

Every packet in $\mathcal Q_{R,\sigma}$ satisfies the hypotheses of
\Cref{conc:line-moments} with $t=s_R$ and $l=l_R$.  Indeed,
$k_s<i_{j+1}^\sigma\le k_R$ gives
\begin{equation}
 d_s>\frac{\rho_{j+1}}{B^2}\ge d_R,
 \qquad \ell_s<l_R.
 \label{hier:residual-precision}
\end{equation}
Since $s_{T(s)}$ and $s_R$ are in the same block-child interval of width
$\rho_{j+1}$,
\[
 |\theta_s-s_R|
 \le |\theta_s-s_{T(s)}|+|s_{T(s)}-s_R|
 \le C d_s+\rho_{j+1}\le C B^2d_s\le C B^4d_s.
\]
For later use, if a residual packet only lies in the larger path node
$J=v_{i_j}(R)$ and has $\ell_s\le l_R$, then
$|s_{T(s)}-s_R|\le\rho_j\le B\rho_{j+1}<B^3d_s$; thus it belongs to
$\mathcal P_R$ as well.  These facts explain both choices of fixed
collections.

\begin{lemma}[Good paths]
\label{hier:good-paths}
There is a fixed exponent $C_{\mathrm{path}}$, chosen independently of
the mass cutoff below, such that the following construction is valid.
Call $R$ bad at $z$ if any of
$\mathcal E_R(z),\mathcal S_R(z),\mathcal V_{R,\sigma}(z)$,
$0\le\sigma<K$, exceeds $L^{C_{\mathrm{path}}}$.  Otherwise call it
good at $z$.  For every fixed angular translation,
\begin{equation}
 \sum_R\int_{\R^2}h_R(z)\one_{\{R\text{ bad at }z\}}\,\dd z
 \le S e^{-10L}.
 \label{hier:weighted-bad-paths}
\end{equation}
The constant $10$ may be replaced by any prescribed absolute constant
by increasing the fixed exponent.  The measurable set
\begin{equation}
 Z=Z_\vartheta
 =\left\{z\in Z_0:
       \sum_{R\text{ bad at }z}h_R(z)\le e^{-L}\right\}
 \label{hier:good-set}
\end{equation}
is independent of $\zeta$ and satisfies
\begin{equation}
 |Z_0\setminus Z|\le S e^{-9L}
 \le L^Ce^{-5L}|F|.
 \label{hier:good-set-budget}
\end{equation}
Its removal is therefore allowable in the target estimate
\eqref{eq:src-7}, using the original adjoint sum's $L^2$ bound.
\end{lemma}

\begin{proof}
Fix $R$, a line of slope $s_R$, and any horizontal interval $I$ of length
$l_R$ on that line.  By \Cref{conc:line-moments}, the normalized
$L^b(I)$ norm of each indicated quantity is at most
$C b^C(1+\log B)^C$.  The same bound holds for
$\mathcal V_{R,\sigma}$ because
\eqref{hier:deterministic-residual-subset} is a fixed subset satisfying
\eqref{eq:src-9}.  As $\log B=K\log G$ is bounded by a fixed power of
$L$, this bound is polynomial in $L$ when $b$ is a fixed positive
multiple of $L$.

For clarity, fix an arbitrarily large absolute number $J_0$.  Choose
$b=\lceil L\rceil$ and then a fixed exponent
$C_{\mathrm{path}}$ sufficiently larger than those in the moment bound.
Markov's inequality gives, for all large $L$,
\[
 \begin{aligned}
 &|\{x\in I:\mathcal E_R(x)>L^{C_{\mathrm{path}}}\}|
  +|\{x\in I:\mathcal S_R(x)>L^{C_{\mathrm{path}}}\}|\\
 &\qquad
  +\sum_{\sigma=0}^{K-1}
   |\{x\in I:\mathcal V_{R,\sigma}(x)>L^{C_{\mathrm{path}}}\}|
  \le e^{-J_0L}|I|.
 \end{aligned}
\]
The factor $K+2$ is polynomial and is absorbed by the same choice.
Alternatively one may increase the fixed multiple in the choice of
$b$; neither choice depends on the number of scales or tops.

Partition the whole reference line into intervals of length $l_R$.
Along this line $h_R$ changes by at most an absolute factor on each such
interval, directly from \eqref{eq:src-5}.  Hence the preceding estimate
implies
\[
 \int_I h_R\one_{\{R\text{ bad}\}}\,\dd x
 \le C e^{-J_0L}\int_I h_R\,\dd x.
\]
Sum over the intervals and integrate in the transverse coordinate
$y-s_Rx$; this change of coordinates has determinant one.  Using
\eqref{hier:top-weight-basic} gives
\[
 \int_{\R^2}h_R\one_{\{R\text{ bad}\}}\,\dd z
 \le L^C e^{-J_0L}|R_R|.
\]
Sum over $R$ and take $J_0$ large enough to obtain
\eqref{hier:weighted-bad-paths}.  This establishes the estimate with
its full spatial weight, not only on the central part of $R$.

Markov's inequality applied to the nonnegative sum defining
\eqref{hier:good-set} gives its first measure bound.  Since
$S\le L^C|E|=L^CN^2|F|$ and $N\le e^{2L}$, the second follows.
For example, Cauchy--Schwarz and the previously proved adjoint estimate
give
\[
 \int_{Z_0\setminus Z}\norm{P}
 \le |Z_0\setminus Z|^{1/2}\norm{P}_2
 \le L^Ce^{-5L/2}N|F|,
\]
which is smaller than any required fixed negative power of $L$ times
$N|F|$.  The estimate is applied to the original sum $P$, as required
for exceptional discards.

Finally, all packet collections used here depend only on the finite
geometric data, the angular translation, and the shift.  They do not
depend on $z$, $\zeta$, significance, or a mass cutoff.  Their finite
length variations and the other thresholded quantities are measurable
in $z$.  They are measurable jointly with $\vartheta$ as well: at the
finitely many relevant angular levels, membership of each of the
finitely many top slopes in its interval changes at finitely many
translation endpoints.  On the complementary parameter intervals
all the relevant membership data are fixed.  This also proves the
asserted independence from $\zeta$.
\end{proof}

\subsection{Mass stopping and the exact path expansion}

Choose a cutoff
\begin{equation}
 M=L^{A_M},
 \label{hier:mass-cutoff}
\end{equation}
where the fixed exponent $A_M$ will be taken sufficiently large in the
residual estimates.  Every threshold and exponent in the preceding
lemmas has already been fixed without reference to $A_M$.  Statements
below that use $M$ hold for each such fixed exponent after increasing
the lower threshold for $L$.

For convenience set
\begin{equation}
 C_M=L^{Cp}(1+M)^{\alpha-1},
 \label{hier:stop-cost}
\end{equation}
with a fixed sufficiently large $C$.  Equation~\eqref{eq:src-17} and
$p\mu=\alpha-1$ imply that a node of mass $n_J\le M$ satisfies
\begin{equation}
 V_J\le C_M n_J.
 \label{hier:small-node-amount}
\end{equation}
Indeed $V_J/n_J\le1+L^{Cp}n_J^{\alpha-1}$ for the parameters.

\begin{lemma}[Expansion into top paths]
\label{hier:path-expansion}
Fix $z\in Z$, a root, a shift, and parameters as above.  Starting
from a node with $n_J>M$, expand
\begin{equation}
 V_J=\beta_J
      \left(n_J-\sum_c n_c+\sum_c V_c\right),
 \label{hier:exact-expansion}
\end{equation}
and stop a branch on first reaching a node of mass at most $M$.
At each expanding node, split $n_J-\sum_c n_c$ into the individual
weights of the tops which do not remain eligible at the block endpoint.
At a stopped node split $V_J$ proportionally among its eligible top
weights.  This gives exactly one terminal path for each eligible root
top $R$, with a terminal amount $e_R$ satisfying
\begin{equation}
 0<e_R\le C_Mh_R,
 \qquad
 V_{\mathrm{root}}
 =\sum_R e_R\prod_{J\text{ active on }R\text{'s path}}\beta_J.
 \label{hier:top-path-expansion}
\end{equation}
If $n_{\mathrm{root}}>M$, at least one half of this sum comes from tops
which are good at $z$.  Consequently,
\begin{equation}
 V_{\mathrm{root}}(z)
 \le 2C_M\sum_{\substack{R\text{ eligible at root}\\R\text{ good at }z}}
     h_R(z)\prod_{J\text{ active on }R\text{'s path}}\beta_J(z).
 \label{hier:good-path-domination}
\end{equation}
If $n_{\mathrm{root}}\le M$, \eqref{hier:small-node-amount} applies
instead.
\end{lemma}

\begin{proof}
Identity~\eqref{hier:exact-expansion} follows directly from
\eqref{eq:src-19}, because
$V_c-n_c=n_c^{\alpha-p}\norm{A_c}^p$.  The term
$n_J-\sum_c n_c$ is a sum of top weights by
\eqref{hier:mass-loss}, now between the endpoints of the whole block.
Every top either remains eligible in exactly one block child or loses
eligibility and terminates at $J$.  The construction therefore branches by disjoint
top sets, assigning each top exactly one path.  It terminates after
finitely many blocks even if no mass stop occurs, since no top is
eligible beyond the final level.  At such a final expanding node the
block-child sum is empty and all tops terminate directly.

If $R$ terminates directly at an expanding node, its terminal amount is
$h_R$.  If it reaches a mass-stopped node $c$, the terminal amount is
$(V_c/n_c)h_R$, which is at most $C_Mh_R$ by
\eqref{hier:small-node-amount}.  Repeated use of
\eqref{hier:exact-expansion} proves the equality in
\eqref{hier:top-path-expansion}.

To verify the claim about bad paths, it is useful to normalize the
expansion by actual probabilities.  At an active node define
\[
 D_J=n_J-\sum_c n_c+\sum_cV_c
     =n_J+\sum_c(V_c-n_c)>M.
\]
Choose a block child $c$ with probability $V_c/D_J$, or a directly
terminating top $R$ with probability $h_R/D_J$.  These probabilities sum
to one.  At a mass-stopped node choose an eligible top with probability
$h_R/n_c$.  The probability of a given complete terminal path is exactly
its summand in \eqref{hier:top-path-expansion} divided by
$V_{\mathrm{root}}$.  In fact, for an active edge $J\to c$,
$V_c/D_J=\beta_JV_c/V_J$, and at a direct terminal edge
$h_R/D_J=\beta_Jh_R/V_J$.  The successive $V_c$ factors cancel;
at a mass stop the last factor $h_R/n_c$ leaves the amount
$(V_c/n_c)h_R$.  This proves the asserted probability identity without
an independence assumption on the amplitudes.

For a directly terminating top, its last conditional probability is
$h_R/D_J\le h_R/M\le h_R$.  If a top terminates in the next node by a
mass stop, the combined last choice of that child and then that top has
conditional probability
\[
 \frac{V_c}{D_J}\frac{h_R}{n_c}
 \le C_M\frac{h_R}{M}\le C_Mh_R.
\]
Earlier conditional probabilities are at most one.  Thus the
unconditional probability of any specified terminal top is at most
$C_Mh_R$.  This estimate also covers tops which terminate after many
active ancestors; it uses only the last active step.  By
\eqref{hier:good-set}, the total probability of bad terminal tops is
at most
\[
 C_M\sum_{R\text{ bad at }z}h_R(z)\le C_Me^{-L}<\frac12
\]
for all sufficiently large $L$.  Here
$\log C_M=O(p\log L+\log(1+M))=o(L)$ for each fixed $A_M$.
The probability identity now implies that good tops carry at least half
of the total amount, and the bound on $e_R$ proves
\eqref{hier:good-path-domination}.
\end{proof}

The remaining path estimate, proved below in
\Cref{res:path-product}, is
\begin{equation}
 \int_{Z\cap\{R\text{ good}\}}
       h_R\prod_{J\text{ active on }R\text{'s path}}\beta_J
 \le L^{Cp}(1+M)^C\int_{\R^2}h_R.
 \label{eq:src-22}
\end{equation}
\Cref{pred:section,res:section} prove \eqref{eq:src-22} from the
estimates established above.  Once that estimate is available,
\eqref{hier:good-path-domination} and
\eqref{hier:top-weight-basic} control the integral of
$V_{\mathrm{root}}$ by summing over the eligible root tops.

\subsection{Path drops and marked spatial cells}

Fix a top $R$ and one shift.  Let $J_R$ be the last block index with
$i_{J_R}\le k_R$.  For $0\le j\le J_R$ put
\begin{equation}
 J_j=v_{i_j}(R),\qquad n_j=n_{J_j},\qquad
 n_{J_R+1}=0,\qquad
 h_j^*=n_j-n_{j+1},\qquad u_j=\frac{h_j^*}{n_j}.
 \label{hier:path-drops}
\end{equation}
For $j<J_R$, the path block child is $v_{i_{j+1}}(R)$.
For $j=J_R$ there is no path child: $R$ ceases to be eligible at the
next block endpoint.  Other tops in $J_j$ may still continue, and all
their children are retained in \eqref{eq:src-18} and
\eqref{eq:src-19}.  Only the path notation sets the next mass to zero.
In particular, the last possible path drop has $h_j^*=n_j$ and $u_j=1$.

The actual active indices form the prefix before the first mass at most
$M$, equivalently the indices with $n_j>M$.  A path which stops by mass
need not reach its final possible block $J_R$.  At the last active
index of such a path we retain the actual eligible next mass in
\eqref{hier:path-drops}; it is not set to zero merely because the
expansion stops there.

\begin{lemma}[Drops and marking]
\label{hier:marking}
At a point of $Z_0$, the active path drops satisfy
\begin{equation}
 \sum_{j:\,n_j>M}u_j
 \le 1+\log^+\frac{n_0}{M}\le C L.
 \label{hier:drop-budget}
\end{equation}
On each line of slope $s_R$, mark an $H_j$ cell $I$ when
\begin{equation}
 \inf_{z\in I}n_j(z)>M,
 \qquad
 \sup_{z\in I}h_j^*(z)\ge\frac12.
 \label{eq:src-23}
\end{equation}
Every point of a marked cell is active and has $h_j^*\ge c$ for an
absolute $c>0$.  At an active point in an unmarked cell with
$h_j^*\ge1$, one has $M<n_j<M+1$ for all sufficiently large $L$.
There is at most one such index on each active path, and its factor is
bounded at points of $Z_0$ by
\begin{equation}
 \beta_{J_j}\le L^{Cp}(1+M)^C.
 \label{hier:missing-factor}
\end{equation}
If an index is inactive at a point, the cell containing that point is
unmarked.  Finally, even in the last possible block,
\begin{equation}
 H_j\le\frac{l_R}{B\sqrt G}.
 \label{hier:mesh-within-top}
\end{equation}
\end{lemma}

\begin{proof}
If there are no active indices there is nothing to prove in
\eqref{hier:drop-budget}.  Otherwise let $q$ be the last active index.
For $j<q$, the next mass is also greater than $M$, and
\[
 u_j=1-\frac{n_{j+1}}{n_j}
 \le\log\frac{n_j}{n_{j+1}}.
\]
The last term $u_q$ is at most one, whether its next mass is positive
or zero.  Therefore the sum is at most
$1+\log(n_0/n_q)\le1+\log(n_0/M)$.  This argument deliberately does
not telescope through a possibly tiny mass below the stopping cutoff.
Equation~\eqref{eq:src-12} gives the last inequality.

Both $n_j$ and $h_j^*$ are positive sums of weights of tops eligible
at $J_j$, except that the difference can be identically zero.  In the
last possible block it is simply $n_j$.  Thus
\Cref{hier:mass-motion} gives relative variation
$1+O(G^{-1/2})$ across the whole $H_j$ cell.  A marked cell has positive
supremum of $h_j^*$, so this difference is not identically zero there,
and its infimum is at least
$\tfrac12 e^{-CG^{-1/2}}\ge\tfrac14$ for large $L$.
The first marking condition supplies activity everywhere on the cell.

Suppose instead that $z$ is active, its cell is unmarked, and
$h_j^*(z)\ge1$.  The second marking condition then holds, so the first
one fails: the cell has infimum mass at most $M$.  The relative-motion
bound implies
\[
 M<n_j(z)\le e^{CG^{-1/2}}M<M+1,
\]
where the last inequality holds because $M$ is a fixed polynomial in
$L$.  Its drop of at least one gives $n_{j+1}(z)<M$ (or zero at the
last possible block), and all subsequent masses are at most that mass.
This is consequently the last active index and can occur at most once.
At a point of $Z_0$, since the denominator of \eqref{eq:src-19} is at
least $n_j$,
\eqref{eq:src-17} gives
\[
 \beta_{J_j}
 \le\frac{V_{J_j}}{n_j}
 \le1+L^{Cp}n_j^{\alpha-1}
 \le L^{Cp}(1+M)^C,
\]
proving \eqref{hier:missing-factor}.  If $n_j(z)\le M$, its cell
cannot have infimum greater than $M$, proving the assertion about
inactive indices.  This observation will justify assigning the value
one to the auxiliary phase factor on every unmarked cell.

Finally, eligibility at the starting level gives
$d_R\le\rho_j/B^2$, equivalently $l_R\ge B^2w/\rho_j$.
Together with \eqref{hier:block-ratio}, this yields
\[
 \frac{H_j}{l_R}
 \le\frac{\rho_j}{B^2\rho_{j+1}\sqrt G}
 \le\frac1{B\sqrt G}.
\]
This only uses eligibility at the start of the block, and therefore
applies to the final possible block as well.
\end{proof}

All constructions above are finite and measurable.  Their constants are
uniform in the finite depth, so no lower bound on a nonzero top weight,
or on the number of angular levels retained, is part of the argument.
The limits from finite packet collections and the angular-translation
average are taken only after the subsequent sections have obtained the
uniform estimates.
